\chapter{Kết luận và hướng phát triển}
\label{chap:conclusion}

\section{Kết luận}

Báo cáo nghiên cứu bài toán nhận dạng ảnh trên dữ liệu đuôi-dài, nơi một số ít lớp phổ biến áp đảo
rất nhiều lớp hiếm (đuôi chỉ 5 ảnh/lớp), theo một mạch tiếp cận có cấu trúc: \emph{đo cho đúng
$\rightarrow$ sửa bộ phân loại $\rightarrow$ mượn tri thức ngoài}. Trên hai bộ dữ liệu khác biệt cả
về quy mô lẫn miền --- CIFAR-100-LT (IF$=$100) và CUB-200-LT (fine-grained, IF$=$10) --- nhóm rút ra
một kết luận nhất quán: với đuôi ít ảnh, chìa khoá \emph{không} nằm ở việc huấn luyện một mạng lớn
hơn từ đầu, mà ở việc \textbf{mượn và thích nghi nhẹ một mô hình nền tảng đóng băng bằng loss
nhận-biết-đuôi}. Trong khi track from-scratch chỉ chạm trần khoảng $0.47$ balanced accuracy (CIFAR)
và $0.20$--$0.27$ (CUB), track nền tảng đạt tới $\mathbf{0.82}$ / $\mathbf{0.85}$. Xét riêng các
nguồn tri thức ngoài, \textbf{mô hình nền tảng thị giác thứ hai (DINOv2) mang lại lợi ích lớn nhất};
ngôn ngữ (LLM) giúp được ít, còn sinh dữ liệu và mixup gần như không giúp.

\section{Những đóng góp chính của nhóm}

\begin{itemize}
  \item \textbf{Xây dựng} một khung thực nghiệm thống nhất, đa-dataset cho long-tailed recognition,
        với giao thức chọn-mô-hình-trên-validation nghiêm ngặt và bộ độ đo tail-aware (balanced
        accuracy, few-shot accuracy, G-mean).
  \item \textbf{Triển khai và đối chứng} đầy đủ ba lớp phương pháp (huấn luyện từ đầu, tái dùng
        checkpoint, thích nghi mô hình nền tảng đóng băng) trên cùng một thước đo.
  \item \textbf{Thực hiện} một nghiên cứu so sánh có hệ thống ba nguồn tri thức ngoài (ngôn ngữ /
        thị giác thứ hai / sinh dữ liệu) trên cùng một đuôi, kèm công cụ đo công bằng (GLA) và phép
        kiểm bù trừ (fusion nhận-biết-đuôi), có đối chiếu với mốc from-scratch.
  \item \textbf{Chứng minh bằng số} (qua hai ablation) rằng phần lớn giá trị đến từ \emph{pipeline
        tail-aware} của nhóm chứ không chỉ từ backbone, và rằng \emph{fine-tune nặng làm hại đuôi}
        --- qua đó biện minh cho thiết kế đóng băng.
  \item \textbf{Báo cáo trung thực} cả các kết quả âm lẫn các ghi chú thực nghiệm, bảo đảm tính khách
        quan.
\end{itemize}

\section{Hạn chế}

Đề tài còn một số hạn chế cần nhìn nhận thẳng thắn. Trước hết là về backbone: do ràng buộc GPU
Kaggle, nhóm chỉ dùng CLIP ViT-B/32 (cấu hình nhẹ nhất), nên track ngôn ngữ (CLIP-LLM) có thể bị
đánh giá thấp --- một CLIP lớn hơn (ViT-L/14) hoàn toàn có thể làm thay đổi cán cân giữa ngôn ngữ và
thị giác. Tiếp đến là vấn đề validation thiếu lớp đuôi: vì lớp tail chỉ có 5 ảnh và được giữ nguyên
cho train, việc tinh chỉnh trên validation không có tín hiệu trực tiếp cho nhóm few, buộc trọng số
nhóm này phải lấy theo medium. Ngoài ra, IF của CUB bị giới hạn ở mức $\le$50 do pool chỉ khoảng 50
ảnh/lớp, nên không gắt bằng CIFAR và được trình bày như một miền fine-grained. Cuối cùng là hai
caveat về breakdown đã nêu trong phần thảo luận (Chương~\ref{chap:results}) liên quan tới bộ
checkpoint Phase 0 và ngưỡng nhóm-shot của CUB from-scratch; cả hai chỉ ảnh hưởng tới breakdown chứ
không ảnh hưởng balanced accuracy.

\section{Hướng phát triển}

Từ những hạn chế trên, một số hướng mở rộng tự nhiên có thể kể tới. Hướng đầu tiên là thử các
backbone nền tảng lớn hơn (CLIP ViT-L/14, DINOv2 ViT-L/14) để kiểm tra xem kết luận ``DINOv2 $>$
CLIP'' có còn giữ nguyên khi tăng dung lượng mô hình hay không. Hướng thứ hai là chạy lại một đối
chứng from-scratch \emph{nhất quán hoàn toàn} --- Phase 0 trên đúng checkpoint
\texttt{pretrained=False}, và from-scratch CUB với ngưỡng nhóm-shot $20/10$ --- nhằm thu được một
bảng ``sạch'' tuyệt đối. Bên cạnh đó, có thể nghiên cứu cách \emph{ước lượng phân phối đuôi} (hoặc
dựng một validation tổng hợp cho đuôi) để chọn siêu tham số fusion cho nhóm few mà không phải buộc
theo medium. Sau cùng, nhóm muốn khảo sát các cơ chế bù trừ sâu hơn giữa các nguồn tri thức (ví dụ
kết hợp ở mức feature thay vì mức xác suất) và mở rộng nghiên cứu sang các dataset long-tail quy mô
lớn hơn như ImageNet-LT hay iNaturalist.
